\chapter{Состояния, обратимость и кинематика}

Классическая механика — это не только набор уравнений для падающих тел, обращающихся планет или заряженных частиц. Прежде чем выбирать какое-либо конкретное уравнение, можно задать более общий вопрос: каким должен быть закон движения, чтобы считаться классическим законом? Мы начнем с систем настолько простых, что любую возможную историю можно нарисовать на доске. Когда идеи состояния, детерминизма, обратимости и сохранения станут там ясны, мы вернемся к частицам, движущимся в обычном пространстве.

\section{Два вида вопросов}

В механике есть два вида вопросов. Один из них касается конкретного закона, управляющего данной системой. Планета, движущаяся в гравитационном поле массивного тела, подчиняется одному закону; электрически заряженная частица, движущаяся в магнитном поле, — другому. Это вопросы о детальных законах природы.

Другой вопрос касается структуры, общей для всех таких примеров. Существуют ли правила, которым должен удовлетворять всякий допустимый закон? Какие черты принадлежат самой классической механике, независимо от силы или объекта, который мы изучаем? Конкретные примеры проиллюстрируют ответ, но именно этот второй вопрос и организует предмет.

Самый полезный первый пример — абстрактная монета на абстрактном столе. Мы сохраняем только одно свойство монеты: она может показывать орла или решку. Поэтому ее пространство конфигураций равно
\begin{equation}
  \mathcal C=\{H,T\}.
\end{equation}
Слово \emph{конфигурация} означает один из возможных способов, которыми может быть устроена система. Начальное условие — это выбор одного из них. Односторонняя монета была бы еще проще, но у нее была бы только одна конфигурация, и ничего никогда не могло бы произойти. Двух состояний достаточно, чтобы динамика стала возможной.

\section{Монета: состояние, закон и история}

Пусть время идет целыми шагами, как если бы мир освещался стробоскопом:
\begin{equation}
  t=\ldots,-2,-1,0,1,2,\ldots .
\end{equation}
Закон движения — это правило обновления. Оно сопоставляет состоянию в момент \(t\) состояние в момент \(t+1\). Самое простое правило ничего не меняет:
\begin{equation}
  H\longmapsto H,
  \qquad
  T\longmapsto T.
  \label{eq:coin-identity}
\end{equation}
Если начать с \(H\), вся история имеет вид \(H,H,H,\ldots\); если начать с \(T\), она имеет вид \(T,T,T,\ldots\). Это правило скучно, но полностью определяет эволюцию.

Другое обратимое правило переставляет два состояния местами:
\begin{equation}
  H\longmapsto T,
  \qquad
  T\longmapsto H.
  \label{eq:coin-flip}
\end{equation}
Теперь истории чередуются:
\begin{equation}
  H,T,H,T,\ldots,
  \qquad
  T,H,T,H,\ldots .
\end{equation}
Стрелка указывает из состояния до обновления в состояние после него. Поэтому хвост стрелки показывает, с чего система начинает, а ее наконечник — куда она переходит.

\begin{figure}[htbp]
  \centering
  \includegraphics[width=0.86\textwidth]{lecture_01_figure_02.png}
  \caption{Две конфигурации монеты и ориентированный закон переворота на лекционной доске. Четкое правило обновления приведено в уравнении~\eqref{eq:coin-flip}.}
  \lectureframe{00:06:50}
  \label{fig:coin-board}
\end{figure}

Стороны можно закодировать двузначной переменной,
\begin{equation}
  \sigma=
  \begin{cases}
    +1,&H,\\
    -1,&T.
  \end{cases}
\end{equation}
Буква \(\sigma\) традиционно используется в физике для двузначных величин, особенно в описаниях спина, хотя никакого знания о спине здесь не требуется. В этих обозначениях два правила становятся уравнениями движения:
\begin{align}
  \sigma(t+1)&=\sigma(t),
  &&\text{identity law},\label{eq:sigma-identity}\\
  \sigma(t+1)&=-\sigma(t),
  &&\text{flip law}.\label{eq:sigma-flip}
\end{align}
Сама переменная не добавляет новой физики. Она дает компактную метку для конфигурации и превращает граф в уравнение.

\begin{figure}[htbp]
  \centering
  \includegraphics[width=0.86\textwidth]{lecture_01_figure_03.png}
  \caption{Закон переворота, записанный как
  \(\sigma(t+1)=-\sigma(t)\). Заключенные в рамку значения \(\sigma=\pm1\) обозначают орла и решку.}
  \lectureframe{00:08:58}
  \label{fig:sigma-board}
\end{figure}

Для любого из двух правил начальное состояние и закон определяют состояние в любой более поздний момент времени. Это утверждение предполагает, что система \emph{замкнута}. Если кто-то протянет руку и перевернет монету, то эта рука будет внешним воздействием, не включенным в наше описание. Мы не опровергли детерминизм; мы изменили систему. Классический закон предназначен управлять замкнутой системой, в описание которой включены все воздействия, существенные для ее эволюции.

Существует 4 математические функции из двухэлементного множества в себя. Те две, которые выше не записаны, сопоставляют обоим входам \(H\) либо обоим входам \(T\). Они однозначно задают будущее, но стирают различие между двумя возможными прошлыми. Когда ниже будет наложено требование обратимости, уравнения~\eqref{eq:coin-identity} и \eqref{eq:coin-flip} окажутся в точности двумя допустимыми отображениями.

\section{Игральная кость: орбиты и сохраняющиеся величины}

Заменим монету игральной костью. Для наших целей важен лишь один факт об игральной кости: у нее 6 конфигураций,
\begin{equation}
  \mathcal C=\{1,2,3,4,5,6\}.
\end{equation}
Начальное условие — это выбор одного числа. Мы снова могли бы оставить каждое состояние неподвижным, но ориентированные графы полезнее для менее тривиальных правил. Рассмотрим
\begin{equation}
  1\longmapsto2\longmapsto3\longmapsto4
  \longmapsto5\longmapsto6\longmapsto1.
  \label{eq:six-cycle}
\end{equation}
Если, например, начать с \(4\), то история будет такой

\[
  4,5,6,1,2,3,4,\ldots .
\]
Шесть состояний образуют один цикл.\footnote{Примечание к источнику: при первом чтении стрелок в одной устной фразе говорится, что \(6\) возвращается в \(6\). Сразу следующая фраза, нарисованный цикл и перечисленная история все указывают на подразумеваемый переход \(6\mapsto1\).}

\begin{figure}[htbp]
  \centering
  \includegraphics[width=0.86\textwidth]{lecture_01_figure_06.png}
  \caption{Единственный цикл из 6 состояний. При каждом стробоскопическом шаге состояние переходит в следующее и возвращается после 6 обновлений.}
  \lectureframe{00:12:40}
  \label{fig:six-cycle-board}
\end{figure}

Правило может выглядеть сложнее, не будучи структурно иным. Например,
\begin{equation}
  1\to2\to5\to3\to4\to6\to1
  \label{eq:relabeled-cycle}
\end{equation}
это по-прежнему один цикл длины 6. Если расположить вершины на странице в порядке \(1,2,5,3,4,6\), получится та же схема, что и в уравнении~\eqref{eq:six-cycle}. Метки вроде \(1\) и \(5\) — это имена; структура цикла от них не зависит.

По-настоящему другой закон разбивает 6 состояний на 2 цикла, например
\begin{equation}
  1\to2\to3\to1,
  \qquad
  4\to6\to5\to4.
  \label{eq:two-cycles}
\end{equation}
Игральная кость по-прежнему одна и закон один, но траектория, которая начинается на одном цикле, никогда не может попасть на другой. Можно приписать циклам число \(Q\),
\begin{equation}
  Q=
  \begin{cases}
    0,&\text{on the cycle }\{1,2,3\},\\
    1,&\text{on the cycle }\{4,5,6\}.
  \end{cases}
\end{equation}
Состояние меняется на каждом шаге, а \(Q\) — нет.

\begin{definition}
Сохраняющаяся величина — это функция состояния, значение которой не изменяется под действием закона движения.
\end{definition}

Та же идея допускает одну неподвижную точку, один двухцикл и один трехцикл. Сохраняющаяся метка указывает, на какой орбите находится система; она не обязана определять состояние внутри этой орбиты. Таким образом, сохранение совместимо с движением. Сохраняется не обязательно полная конфигурация, а некоторое различие между возможными конфигурациями.

\begin{figure}[htbp]
  \centering
  \includegraphics[width=0.86\textwidth]{lecture_01_figure_07.png}
  \caption{Несколько несвязанных циклов на доске: неподвижная точка, двухцикл и трехцикл. Метка цикла сохраняется даже тогда, когда состояние движется по своему циклу.}
  \lectureframe{00:18:40}
  \label{fig:conserved-cycles-board}
\end{figure}

В конечной обратимой системе всякая траектория является циклом. У единственного цикла нет непостоянной сохраняющейся метки такого элементарного вида, потому что траектория рано или поздно проходит через каждое состояние. Несколько циклов допускают такую метку, поскольку эволюция никогда не переносит состояние из одной компоненты в другую.

\section{Что должно входить в состояние?}

В примерах с монетой использовалось одношаговое правило: следующая сторона зависела только от текущей стороны. Вопрос из аудитории проверяет, существенно ли это ограничение.

\begin{classroomqa}
\audiencequestion{Можно ли построить последовательность орлов и решек, в которой следующий результат зависит не только от непосредственно видимой стороны, а, возможно, от одного или нескольких предыдущих результатов?\lecturetimestamp{00:19:27--00:21:23}}

\lecturerresponse{Да. В этом случае одних \(H\) или \(T\) недостаточно для полного задания состояния. Если нужна одна предыдущая сторона, расширьте пространство состояний до
\(
  \{HH,HT,TH,TT\}.
\)
Каждая упорядоченная пара записывает текущую сторону и предыдущую. Тогда эволюцию снова можно записать как одношаговое правило на этом пространстве из четырех состояний. Если требуется еще больше предыстории, включите и эту информацию.}
\end{classroomqa}

Этот ответ дает операциональный смысл состояния.

\begin{definition}
Состояние системы в некоторый момент состоит из всей информации, необходимой вместе с законом движения для предсказания ее будущего.
\end{definition}

Следовательно, состояние — не просто фотография. Это любой полный набор данных, делающий эволюцию однозначной. Позже, когда мы будем описывать ньютоновскую частицу, одного ее положения будет недостаточно; к начальным данным будет относиться и скорость. Монета с памятью — простейший вариант того же принципа.

\section{Бесконечные пространства состояний}

Бесконечное число состояний не требует бесконечного числа объектов. Рассмотрим одну частицу, ограниченную целочисленными узлами бесконечной прямой. Ее состояние задается целым числом
\[
  n\in\mathbb Z.
\]
Один закон сдвигает ее на один узел вправо на каждом шаге:
\begin{equation}
  n(t+1)=n(t)+1.
  \label{eq:lattice-one}
\end{equation}
Эта орбита никогда не замыкается, но проходит через каждое целое число, если проследить ее как в будущее, так и в прошлое. В этом расширенном смысле она играет роль одного бесконечного цикла.

Другой закон сдвигает ее сразу на два узла:
\begin{equation}
  n(t+1)=n(t)+2.
  \label{eq:lattice-two}
\end{equation}
Теперь четные узлы переходят только в четные, а нечетные — только в нечетные. Четность сохраняется:
\begin{equation}
  \begin{aligned}
    Q_{\rm parity}(n)&=n\bmod 2,\\
    Q_{\rm parity}(n(t+1))&=Q_{\rm parity}(n(t)).
  \end{aligned}
\end{equation}
Переход от конечного пространства конфигураций к бесконечному допускает истории, которые никогда не повторяются, но не меняет смысла понятий состояния, орбиты или сохранения.

\section{Детерминизм — это еще не обратимость}

Теперь введем правило, которое однозначно предсказывает будущее, но не удовлетворяет второму классическому требованию. Припишем трехсторонней монете состояния орел, решка и ребро и определим
\begin{equation}
  H\to T,
  \qquad
  T\to E,
  \qquad
  E\to T.
  \label{eq:nonreversible-map}
\end{equation}
Начиная с \(H\), история имеет вид
\[
  H,T,E,T,E,\ldots;
\]
а начиная либо с \(T\), либо с \(E\), последующее чередование столь же однозначно. Из каждого состояния выходит ровно одна стрелка, так что это правило детерминистично по отношению к будущему.

Трудность возникает, когда мы пытаемся восстановить прошлое. Если текущее состояние — \(T\), то предыдущим могло быть \(H\) или \(E\). Две истории слились. Если обратить каждую стрелку, изъян сразу становится очевидным: из \(T\) на обращенном графе возможны два хода. Обратное правило не является детерминистическим.

\begin{figure}[htbp]
  \centering
  \includegraphics[width=0.86\textwidth]{lecture_01_figure_08.png}
  \caption{Необратимый закон для \(H,T,E\) и характерные истории. И \(H\), и \(E\) переходят в \(T\), так что настоящее не определяет единственного прошлого.}
  \lectureframe{00:28:50}
  \label{fig:nonreversible-board}
\end{figure}

\begin{definition}
Дискретный закон обратим, если текущее состояние определяет как единственное следующее состояние, так и единственное предыдущее состояние. В графе состояний из каждого состояния должна выходить ровно одна стрелка и в него должна входить ровно одна стрелка.
\end{definition}

Для конечного пространства состояний обратимое обновление представляет собой перестановку состояний. Оно никогда не создает и не уничтожает различий между состояниями; оно лишь переставляет их. Поэтому его граф распадается на циклы. Отображение в уравнении~\eqref{eq:nonreversible-map} является функцией, но не взаимно однозначной функцией, и потому не имеет обратного отображения.

\begin{classroomqa}
\audiencequestion{Могло бы правило такого рода быть допустимым, если бы состояния представляли вероятности для частицы?\lecturetimestamp{00:29:46--00:31:17}}

\lecturerresponse{Это была бы другая схема. Фундаментальная классическая механика в той формулировке, которую мы здесь используем, опирается на детерминистические законы, а не на вероятностные правила обновления. Конечно, такое правило можно придумать и изучать, но в этом смысле оно не является законом классической механики. Квантовая механика не детерминистична в том же смысле и тем не менее имеет аналог обратимой эволюции; это требует отдельного, более позднего обсуждения. Классическая статистическая механика тоже использует вероятности, но она исходит из лежащих в основе механических законов и описывает неполное знание о начальных данных.}
\end{classroomqa}

Это различие важно. Распределение вероятностей, используемое потому, что мы не знаем точного классического состояния, не является индетерминистическим законом движения. Оно описывает нашу информацию о системе, микроскопическая эволюция которой по-прежнему считается детерминистической и обратимой.

\section{Точные законы и неполное знание}

Для предсказания требуются две составляющие:
\begin{enumerate}
  \item закон движения;
  \item начальное состояние.
\end{enumerate}
Точный закон не дает точного прогноза, если начальное условие известно лишь приблизительно. Для монеты легко представить, что мы смотрим на стол и совершенно точно узнаем, какой стороной она лежит. Реальные механические переменные непрерывны. Положения и скорости могут принимать произвольные вещественные значения, и никакое измерение не дает бесконечно много десятичных знаков.

Предположим, что истинное начальное состояние лежит в небольшой области возможных состояний. При одной динамике эта область долго остается небольшой; при другой близкие начальные состояния быстро расходятся. В хаотической системе малые начальные неопределенности со временем вырастают в большие различия. Атмосфера — хорошо известный пример. Более точные измерения могут продлить интервал полезного прогноза, но при любой фиксированной точности измерения найдется достаточно удаленный момент времени, в который прогноз становится ненадежным.

Это практическое ограничение не вносит неоднозначности в уравнения. Если бы начальные условия были известны точно, классический обратимый закон определял бы и будущее, и прошлое. При конечной точности количественный вопрос вместо этого таков: насколько точно должно быть известно начальное состояние, чтобы сделать предсказание на заданный промежуток времени?

Теперь графический критерий приводит к краткому выводу. Закон, представленный одной входящей и одной исходящей стрелкой в каждом состоянии, детерминистичен и обратим. Соответствующие утверждения для частиц будут менее элементарны, но они выражают ту же классическую структуру.

\section{От абстрактных состояний к точкам в пространстве}

Теперь мы переходим от монет и игральных костей к материальным точкам. Материальная точка может двигаться в одном, трех, пяти или любом другом числе измерений, но непосредственный пример дает обычное пространство. Чтобы количественно задать точку, выберите:
\begin{itemize}
  \item начало координат;
  \item три взаимно перпендикулярные оси;
  \item ориентацию этих осей;
  \item единицу длины.
\end{itemize}
Начало координат выбирается условно. Его можно поместить в Пало-Альто или где угодно еще, и физические выводы не должны зависеть от этого выбора. Оси можно обозначить как \(x,y,z\) или как \(x_1,x_2,x_3\). После того как направления \(x\) и \(y\) зафиксированы, ось \(z\) может быть направлена в любое из двух направлений, перпендикулярных им обоим. Мы выбираем правостороннюю конвенцию: если большой палец правой руки направлен вдоль \(x\), а указательный — вдоль \(y\), то средний палец указывает вдоль \(z\).

Чтобы наносить равные отметки вдоль осей, нужна единица длины, например метры, футы, дюймы или световые годы. После того как эти соглашения зафиксированы, каждая точка задается тремя числами,
\begin{equation}
  (x,y,z).
\end{equation}
В точку можно попасть из начала координат, смещаясь вдоль трех осей. Система из многих частиц требует по одной такой тройке для каждой частицы.

\begin{figure}[htbp]
  \centering
  \includegraphics[width=0.86\textwidth]{lecture_01_figure_09.png}
  \caption{Декартовы оси и три координаты точки. Рисунок фиксирует координатное соглашение; сама точка не зависит от того, где мы выбираем начало координат или оси.}
  \lectureframe{00:43:00}
  \label{fig:coordinates-board}
\end{figure}

\section{Векторы и их компоненты}

Вектор — это объект, имеющий величину и направление. Перемещение от начала координат к точке есть вектор, но векторы также описывают скорость, ускорение, электрические поля и многие другие величины. Мы будем обозначать вектор полужирным шрифтом, например \(\mathbf v\), а его модуль — как \(\lvert\mathbf v\rvert\). Модуль неотрицателен и обращается в нуль только для нулевого вектора.

С геометрической точки зрения вектор — это свободная стрелка. Параллельный перенос стрелки без изменения ее длины или направления не изменяет вектор. В конкретном приложении вектор может быть привязан к точке, например когда мы задаем электрическое поле в данной точке, но такая привязка является дополнительной информацией и не входит в само абстрактное понятие вектора.

Относительно декартовых осей вектор имеет компоненты
\begin{equation}
  \mathbf v=(v_x,v_y,v_z).
\end{equation}
Компоненты — это обычные числа со знаком, получаемые проецированием на оси. Они задают и величину, и направление. Трехмерная теорема Пифагора дает
\begin{equation}
  \lvert\mathbf v\rvert^2
  =v_x^2+v_y^2+v_z^2,
  \qquad
  \lvert\mathbf v\rvert
  =\sqrt{v_x^2+v_y^2+v_z^2}.
  \label{eq:vector-magnitude}
\end{equation}
Отношения компонент определяют направление. Отрицательная компонента означает, что вектор отчасти направлен вдоль отрицательного направления соответствующей оси.

\subsection{Умножение на скаляр и сложение}

Умножение на вещественное число изменяет длину вектора и, если это число отрицательно, меняет его направление на противоположное. При \(c>0\) вектор \(c\mathbf A\) направлен так же, как \(\mathbf A\), и имеет модуль \(c\lvert\mathbf A\rvert\). В частности, \(-\mathbf A\) имеет ту же длину, что и \(\mathbf A\), и направлен противоположно.

Есть три эквивалентных наглядных представления сложения векторов. Во-первых, поместите начало \(\mathbf B\) в конец \(\mathbf A\); стрелка от начала \(\mathbf A\) к концу \(\mathbf B\) есть \(\mathbf A+\mathbf B\). Во-вторых, совместите оба начала и достройте параллелограмм; его диагональ и есть сумма. В этой форме наглядно видна коммутативность:
\begin{equation}
  \mathbf A+\mathbf B=\mathbf B+\mathbf A.
\end{equation}
В-третьих, сложите соответствующие компоненты:
\begin{equation}
  C_i=A_i+B_i,
  \qquad i=x,y,z,
  \label{eq:component-addition}
\end{equation}
где \(\mathbf C=\mathbf A+\mathbf B\). Эти правила также определяют произвольные линейные комбинации, такие как \(a\mathbf A+b\mathbf B\).

\begin{figure}[htbp]
  \centering
  \includegraphics[width=0.86\textwidth]{lecture_01_figure_10.png}
  \caption{Треугольное и параллелограммное построения для сложения векторов. Оба дают один и тот же вектор \(\mathbf A+\mathbf B\).}
  \lectureframe{00:53:00}
  \label{fig:vector-addition-board}
\end{figure}

Общей операции деления одного вектора на другой не существует. Однако существуют два стандартных произведения векторов. Векторное произведение мы оставим для более поздней лекции. Здесь нам нужно скалярное произведение.

\subsection{Скалярное произведение}

Спроецируйте \(\mathbf A\) на ось, задаваемую \(\mathbf B\). Если угол между векторами равен \(\theta\), то алгебраическая компонента \(\mathbf A\) вдоль \(\mathbf B\) равна \(\lvert\mathbf A\rvert\cos\theta\). Умножение этой проекции на \(\lvert\mathbf B\rvert\) определяет
\begin{equation}
  \mathbf A\cdot\mathbf B
  =\lvert\mathbf A\rvert\,\lvert\mathbf B\rvert\cos\theta.
  \label{eq:dot-geometric}
\end{equation}
Результат — скаляр, а не еще один вектор. Уравнение~\eqref{eq:dot-geometric} симметрично относительно \(\mathbf A\) и \(\mathbf B\), так что \(\mathbf A\cdot\mathbf B=\mathbf B\cdot\mathbf A\).

\begin{figure}[htbp]
  \centering
  \includegraphics[width=0.86\textwidth]{lecture_01_figure_11.png}
  \caption{Построение скалярного произведения на доске с помощью проекции. В печатной записи это
  \(\mathbf A\cdot\mathbf B=AB\cos\theta\).}
  \lectureframe{00:58:30}
  \label{fig:dot-product-board}
\end{figure}

Знак несет геометрическую информацию. Скалярное произведение положительно при остром угле, отрицательно при тупом и равно нулю для перпендикулярных векторов. Следовательно,
\begin{equation}
  \mathbf A\perp\mathbf B
  \quad\Longleftrightarrow\quad
  \mathbf A\cdot\mathbf B=0
\end{equation}
для ненулевых векторов.

В компонентах,
\begin{equation}
  \mathbf A\cdot\mathbf B
  =A_xB_x+A_yB_y+A_zB_z.
  \label{eq:dot-components}
\end{equation}
эта формула позволяет выполнять вычисления сразу. В частности,
\begin{equation}
  \mathbf A\cdot\mathbf A
  =A_x^2+A_y^2+A_z^2
  =\lvert\mathbf A\rvert^2.
  \label{eq:self-dot}
\end{equation}
Если заданы компоненты двух векторов, то формулы~\eqref{eq:dot-geometric} и
\eqref{eq:dot-components} определяют угол между ними:
\begin{equation}
  \cos\theta
  =\frac{\mathbf A\cdot\mathbf B}
  {\lvert\mathbf A\rvert\,\lvert\mathbf B\rvert}.
\end{equation}

Закон косинусов следует из той же алгебры. Пусть \(\mathbf A\) и
\(\mathbf B\) проведены из одной вершины к двум другим вершинам треугольника.
Третья сторона есть
\(\mathbf C=\mathbf A-\mathbf B\), а не \(\mathbf A+\mathbf B\). Поэтому
\begin{align}
  C^2
  &=(\mathbf A-\mathbf B)\cdot(\mathbf A-\mathbf B)\notag\\
  &=A^2+B^2-2\mathbf A\cdot\mathbf B\notag\\
  &=A^2+B^2-2AB\cos\theta.
  \label{eq:law-of-cosines}
\end{align}

\begin{figure}[htbp]
  \centering
  \includegraphics[width=0.86\textwidth]{lecture_01_figure_04.png}
  \caption{Треугольник, используемый для закона косинусов. Исправление на доске
  показывает, что замыкающая сторона равна
  \(\mathbf C=\mathbf A-\mathbf B\), что приводит к
  формуле~\eqref{eq:law-of-cosines}.}
  \lectureframe{01:04:29}
  \label{fig:law-cosines-board}
\end{figure}

\section{Положение, скорость и ускорение}

Выберем начало координат, и пусть \(\mathbf r\) — радиус-вектор, проведенный от него
к частице. Его компоненты — декартовы координаты:
\begin{equation}
  \mathbf r=(x,y,z),
  \qquad
  \lvert\mathbf r\rvert=\sqrt{x^2+y^2+z^2}.
\end{equation}
Изменение этого выбора меняет \(\mathbf r\); к этому вопросу мы пока вернемся позже. Когда
оно зафиксировано, движение описывается зависящим от времени положением,
\begin{equation}
  \mathbf r(t)=\bigl(x(t),y(t),z(t)\bigr).
\end{equation}
Следовательно, траектория задается тремя функциями времени.

Скорость есть производная положения по времени:
\begin{equation}
  \mathbf v(t)=\frac{d\mathbf r}{dt}
  =\dot{\mathbf r}(t),
  \qquad
  v_i=\frac{dx_i}{dt}=\dot x_i.
  \label{eq:velocity}
\end{equation}
Точка сверху всегда обозначает производную по времени. Величина
\(\lvert\mathbf v\rvert\) есть модуль скорости. В общем случае направления радиус-вектора и
скорости независимы: частица, находящаяся в данной точке, может двигаться в любом
направлении, допускаемом задачей.

Ускорение есть скорость изменения скорости и, следовательно, вторая производная
положения по времени:
\begin{equation}
  \mathbf a(t)=\frac{d\mathbf v}{dt}
  =\dot{\mathbf v}(t)
  =\frac{d^2\mathbf r}{dt^2}
  =\ddot{\mathbf r}(t),
  \qquad
  a_i=\ddot x_i.
  \label{eq:acceleration}
\end{equation}
Ускорение тоже является вектором. Если частица, движущаяся в положительном
направлении \(x\), увеличивает скорость, то \(a_x>0\). Если она замедляется, то \(a_x<0\). То, что
часто называют замедлением, — это просто ускорение, направленное противоположно текущей
скорости. В более общем случае скорость может изменяться за счет изменения своего направления
даже тогда, когда ее величина остается постоянной.

\begin{figure}[htbp]
  \centering
  \includegraphics[width=0.86\textwidth]{lecture_01_figure_05.png}
  \caption{Компоненты скорости, записанные как производные по времени от компонент
  положения. В компактной векторной записи,
  \(\mathbf v=d\mathbf r/dt\).}
  \lectureframe{01:14:24}
  \label{fig:velocity-board}
\end{figure}

\subsection{Движение на прямой с постоянным ускорением}

Для одномерного движения запишем
\begin{equation}
  x(t)=a+bt+ct^2.
  \label{eq:quadratic-motion}
\end{equation}
При \(t=0\) частица находится в \(x=a\). Однократное, а затем двукратное
diфференцирование дает
\begin{align}
  v(t)=\dot x(t)&=b+2ct,\label{eq:line-velocity}\\
  a(t)=\ddot x(t)&=2c.\label{eq:line-acceleration}
\end{align}
Следовательно, \(b\) — это начальная скорость, а \(2c\) — постоянное ускорение. Если
\(c>0\), скорость линейно возрастает; если \(c<0\), она линейно убывает.
Этот квадратичный закон положения — именно та форма, которую принимает движение при постоянном
ускорении, например идеализированное движение падающего тела вблизи поверхности
Земли.

\section{Равномерное движение по окружности}

В качестве последнего примера пусть частица движется в плоскости \(x\)-\(y\) по
единичной окружности. При \(t=0\) она находится на положительной полуоси \(x\), а ее угловая
координата растет равномерно:
\begin{equation}
  \theta(t)=\omega t.
\end{equation}
Постоянная \(\omega\) есть угловая частота. Одному обороту соответствует
увеличение угла на \(2\pi\), поэтому период равен
\begin{equation}
  T=\frac{2\pi}{\omega}.
  \label{eq:circular-period}
\end{equation}
Можно позволить углу возрастать и дальше \(2\pi\); положение повторяется,
хотя непрерывно отслеживаемый угол не сбрасывается.

На единичной окружности,
\begin{equation}
  \mathbf r(t)
  =\bigl(\cos\omega t,\,\sin\omega t\bigr).
  \label{eq:circular-position}
\end{equation}
Дифференцирование дает
\begin{equation}
  \mathbf v(t)
  =\omega\bigl(-\sin\omega t,\,\cos\omega t\bigr).
  \label{eq:circular-velocity}
\end{equation}
Скорость направлена по касательной к окружности. Ее перпендикулярность
радиусу можно доказать, не опираясь на рисунок:
\begin{align}
  \mathbf r\cdot\mathbf v
  &=\cos\omega t\,(-\omega\sin\omega t)
    +\sin\omega t\,(\omega\cos\omega t)\notag\\
  &=0.
  \label{eq:r-dot-v}
\end{align}

\begin{figure}[htbp]
  \centering
  \includegraphics[width=0.86\textwidth]{lecture_01_figure_12.png}
  \caption{Компоненты положения и скорости при движении по единичной окружности.
  Сокращение в \(\mathbf r\cdot\mathbf v\) показывает, что скорость
  направлена по касательной к окружности.}
  \lectureframe{01:24:30}
  \label{fig:circular-velocity-board}
\end{figure}

Дифференцируя еще раз,
\begin{equation}
  \mathbf a(t)
  =-\omega^2
  \bigl(\cos\omega t,\,\sin\omega t\bigr)
  =-\omega^2\mathbf r(t).
  \label{eq:centripetal-acceleration}
\end{equation}
Ускорение параллельно радиусу, но направлено в противоположную
сторону, к центру. Это центростремительное ускорение.\footnote{Терминологическое
замечание: в этом месте лекции встречается слово ``центробежное'', но
уравнение и сопровождающее описание указывают на стандартное центростремительное
ускорение, направленное к центру.}

\begin{figure}[H]
  \centering
  \includegraphics[width=0.86\textwidth]{lecture_01_figure_13.png}
  \caption{Покомпонентное вычисление ускорения при движении по окружности. Каждая
  компонента равна соответствующей компоненте положения, умноженной на \(-\omega^2\), поэтому
  \(\mathbf a=-\omega^2\mathbf r\).}
  \lectureframe{01:26:30}
  \label{fig:circular-acceleration-board}
\end{figure}

Поскольку радиус равен единице,
\begin{equation}
  \lvert\mathbf r\rvert=1,
  \qquad
  \lvert\mathbf v\rvert=\omega,
  \qquad
  \lvert\mathbf a\rvert=\omega^2.
\end{equation}
При увеличении \(\omega\) модуль скорости растет линейно, а требуемое направленное внутрь
ускорение — квадратично. Для окружности радиуса \(R\) тот же расчет
дал бы \(v=\omega R\) и
\(\lvert\mathbf a\rvert=\omega^2R=v^2/R\); в лекции используется \(R=1\), чтобы не загромождать
первое вычисление.

Теперь у нас есть элементарный язык, необходимый для механики частицы: состояния и
начальные данные, обратимая эволюция, координаты, векторы и производные,
определяющие скорость и ускорение. Следующая задача состоит в том, чтобы точно
определить, какие начальные данные нужно задать для частицы, и найти законы,
которые переносят эти данные от одного момента к следующему.
